\section{One Process, Two Files}
\label{sec:ch10-one-process-two-files}
\index{router!starting it|(}%

The router is a service in \code{docker-compose.yml} like the other three, and
starting it is one command:

\begin{minted}{text}
docker compose up -d mosaic-router
\end{minted}

\index{router!what it needs to start}%
It needs no account, no registry and no network beyond the two subgraphs it
will eventually call. Chapter~\ref{ch:how-a-federated-query-actually-runs}
established that much on the sample; what is new here is that it needs nothing
of Mosaic either. I started it with both services stopped and it came up and
answered \code{/health}, because the execution config is a complete description
of the graph and the router does not check whether any of it is true until a
query arrives. So the compose service has no \code{depends\_on}. There is
nothing to wait for.

Figure~\ref{fig:ch10-what-it-loads} is the whole arrangement: two files in, one
process, two services it reaches by a route that is not quite the one written
down.

\begin{figure}[htbp]
  \centering
  % What the router loads and what it talks to. Referenced from chapter 10.
% Bare tikzpicture; the figure environment, caption and label live at the call
% site.
%
% The point of the picture is the boundary in the middle: the router is in a
% container and the subgraphs are not, which is what makes the routing URLs in
% the composed config wrong and the dashed rewrite necessary.
\begin{tikzpicture}[
    font=\small,
    infile/.style={draw, rounded corners=2pt, minimum width=44mm,
                   minimum height=8mm, font=\scriptsize, align=center},
    proc/.style={draw, rounded corners=2pt, minimum width=62mm,
                 minimum height=12mm, font=\scriptsize, align=center,
                 fill=black!8},
    svc/.style={draw, rounded corners=2pt, minimum width=44mm,
                minimum height=10mm, font=\scriptsize, align=center,
                fill=black!5},
    feed/.style={-{Stealth[length=2mm]}},
    call/.style={-{Stealth[length=2mm]}, dashed, gray},
    lbl/.style={font=\scriptsize, align=center, inner sep=1pt},
    divider/.style={gray, dotted, thick}
  ]

  % -- the two mounted files -------------------------------------------------
  \node[infile] (cfg) at (3.4,0)  {\code{federation/supergraph.json}\\the graph};
  \node[infile] (yml) at (10.4,0) {\code{router/config.yaml}\\the process};

  % -- the router ------------------------------------------------------------
  \node[proc] (router) at (6.9,-2.2)
    {Cosmo Router 0.337.1, in a container\\listening on \code{0.0.0.0:3002}};

  \draw[feed] (cfg.south) -- (3.4,-2.2) -- (router.west);
  \draw[feed] (yml.south) -- (10.4,-2.2) -- (router.east);

  % -- the container boundary ------------------------------------------------
  \draw[divider] (0.3,-3.35) -- (13.5,-3.35);
  \node[lbl, gray, anchor=south east] at (13.5,-3.28) {in a container};
  \node[lbl, gray, anchor=north east] at (13.5,-3.42) {on the host};

  % -- the two subgraphs -----------------------------------------------------
  \node[svc] (cat) at (3.4,-5.5)  {\code{Mosaic.Catalog}\\port 5101};
  \node[svc] (mos) at (10.4,-5.5) {\code{Mosaic.Api}\\port 5100};

  \draw[call] (router.south) -- (6.9,-4.75) -- (3.4,-4.75) -- (cat.north);
  \draw[call] (router.south) -- (6.9,-4.75) -- (10.4,-4.75) -- (mos.north);

  \node[lbl, gray, anchor=south] at (6.9,-4.65)
    {the config says \code{localhost};
     the router dials \code{host.docker.internal}};

\end{tikzpicture}

  \caption{Everything the router knows is in the two mounted files, and neither
  is consulted again after startup unless something changes them. The dashed
  arrow is the part that is not written anywhere: the routing URLs say
  \code{localhost}, and inside a container that is the container, so the router
  rewrites them before it dials. The setting behind that rewrite is in
  section~\ref{sec:ch10-sixty-six-settings}.}
  \label{fig:ch10-what-it-loads}
\end{figure}

\subsection*{What it says on the way up}

\index{router!startup log}%
Eleven lines, of which three are warnings. This is the whole of it, colour
stripped and with the \code{hostname}, \code{pid} and \code{service\_version}
fields that repeat on every line taken out. Nothing else is removed:

\begin{minted}[fontsize=\footnotesize,breaklines]{text}
INFO cmd/main.go:230 Config file watching is disabled, you can still trigger reloads by sending SIGHUP to the router process
INFO cmd/main.go:125 Config file provided. Values in the config file have higher priority than environment variables {"config_file": ["/etc/mosaic-router/config.yaml"]}
INFO core/supervisor_instance.go:47 GOMEMLIMIT set automatically {"limit": "0 B"}
WARN core/router.go:464 No graph token provided. The following Cosmo Cloud features are disabled. Not recommended for Production. {"features": ["Schema Usage Tracking", "Persistent operations", "Cosmo Cloud Tracing", "Cosmo Cloud Metrics"]}
WARN core/router.go:536 Development mode enabled. This should only be used for testing purposes
INFO metric/prometheus_server.go:63 Prometheus metrics enabled {"listen_addr": "127.0.0.1:8088", "endpoint": "/metrics"}
INFO core/router.go:1020 Serving GraphQL playground {"url": "http://0.0.0.0:3002/"}
WARN core/router.go:1570 Advanced Request Tracing (ART) is enabled in development mode but requires a graph token to work in production.
INFO core/router.go:1697 Server initialized and ready to serve requests {"listen_addr": "0.0.0.0:3002", "playground": true, "introspection": true, "config_version": "00000000-0000-0000-0000-000000000000"}
INFO core/router.go:1715 Watching config file for changes. Router will hot-reload automatically without downtime {"path": "/etc/mosaic/supergraph.json"}
INFO core/supervisor.go:155 Router started
\end{minted}

\index{router!config\_version@\code{config\_version}}%
\code{config\_version} is chapter~\ref{ch:composition}'s all-zero \code{version}
field arriving where chapter~\ref{ch:how-a-federated-query-actually-runs} first
met it. A config composed on a laptop has no registry to take a version from,
and the router prints the zeroes without comment.

\index{router!two settings named config file}%
Read the first line and the second to last one together, because they are two
different settings with almost the same name and they say opposite things.
Line one is \code{watch\_config}, which watches \code{config.yaml}, and it is
off. Line ten is \code{execution\_config.file.watch}, which watches
\code{supergraph.json}, and it is on because Mosaic's configuration turns it on.
Nothing on either line says which file it means until the tenth one names a
path, and the two are eight lines apart.

\subsection*{The warnings are the honest kind}

Three warnings on a healthy start would normally be a smell. These three are
each telling you about a decision this book made on purpose.

\index{router!no graph token}%
\code{No graph token provided} is the account this book does not have. The
token is how a router fetches its schema from the Cosmo registry and reports
back to it, and with a static execution config there is nothing to fetch. What
it costs is named in the warning itself: schema usage tracking, persisted
operations, and Cosmo Cloud's tracing and metrics.
Chapter~\ref{ch:self-hosting-cosmo} is where a control plane arrives and those
come back.

\code{Development mode enabled} and the note about Advanced Request Tracing are
the same decision twice. Development mode is what makes the query plan headers
work without a signed request, and
section~\ref{sec:ch10-the-plan-on-a-real-graph} depends on it entirely. It is
a development posture, exactly like the \code{ASPNETCORE\_ENVIRONMENT} of
\code{Development} that both subgraphs run under in the same compose file, and
for the same reason: a graph you cannot introspect is a graph no tool can be
pointed at.

\medskip
\index{router!endpoints}%
Beyond \code{/graphql} the router serves an IDE at its root, and three health
endpoints at their default paths. I asked for all three and got 200 and the
body \code{OK} from each:

\begin{minted}{text}
/health        /health/ready        /health/live
\end{minted}

\code{/metrics} on that port answers 404, because the Prometheus server the log
mentions is on \code{127.0.0.1:8088} inside the container and nothing publishes
it. Chapter~\ref{ch:observability} is where that becomes interesting.

\index{router!readiness reports on itself}%
One thing about readiness will matter later, so here it is measured rather than
assumed. I stopped both subgraphs, started the router against nothing, and
asked all three endpoints again. All three still answer 200 and \code{OK}. A
query at that point comes back
\code{Failed to fetch from Subgraph 'catalog'} with \code{data} null. The
router is reporting on itself, and it is ready in the sense that it is willing
to take your request.

\index{router!starting it|)}%
